\section{Unflavoured S-duality}
\label{sec:unflavoured}
We first consider the setting in which $u=1$.
Throughout this section we will further assume $0<q<1$. Our results will extend
to $0<|q|<1$ by holomorphy of the original index integrals.
In this case the integrands of~\eqref{eq:root-integrals} and~\eqref{eq:root-C}
can be expressed in terms of Jacobi elliptic functions.
More explicitly, take
\begin{equation}
  R(a)=q^{1/2}\sigma(a),\qquad
  \sigma(a)=\frac{\vartheta_1(a)^2}{\vartheta_4(a)^2}
  =k\sn^2\!\left(\frac{2K}{\pi}a\right),
  \label{eq:x}
\end{equation}
with elliptic modulus $k=\vartheta_2(0)^2/\vartheta_3(0)^2$ and complete
elliptic integral $K=\frac{\pi}{2}\vartheta_3(0)^2$. Write
$\sigma_B(a)=\sigma(a)$ and $\sigma_C(a)=\sigma(2a)$.
Substituting~\eqref{eq:x} into~\eqref{eq:root-integrals} and~\eqref{eq:root-C}
and using the addition formula for $\sn$ gives
\begin{equation}
  \mathcal I_G(q,1;n)=\frac{C_0(Q,q)^n q^{-n^2/2}}{2^n n!}
  \int_{[0,\pi]^n}\prod_{i=1}^n\frac{\dd a_i}{\pi}\,
  \prod_{i=1}^n\sigma_G(a_i)
  \prod_{i<j}\left(\frac{\sigma(a_i)-\sigma(a_j)}
  {1-\sigma(a_i)\sigma(a_j)}\right)^2.
  \label{eq:unflavoured-integral}
\end{equation}

\subsection{The Fermi-gas representation}

The integrand can be reduced to a determinant using a variation on Cauchy's
double alternant identity~\cite{Krattenthaler:1999}:\footnote{Take
equation~(2.7) of~\cite{Krattenthaler:1999} and specialize to $X_i=1/x_i$ and
$Y_j=-x_j$; zero values of $x_i$ follow by continuity.}
\begin{equation}
  \det_{i,j}\!\left[\frac{1}{1-x_ix_j}\right]
  =\frac{\prod_{i<j}(x_i-x_j)^2}
  {\prod_i(1-x_i^2)\prod_{i<j}(1-x_ix_j)^2}.
  \label{eq:cauchy}
\end{equation}
With $x_i=\sigma(a_i)$, the integrand in~\eqref{eq:unflavoured-integral} is the
determinant of the symmetric kernel
\begin{equation}
  \widehat\rho_G(a,b)
  =\frac{w_G(a)w_G(b)
    \sqrt{(1-\sigma(a)^2)(1-\sigma(b)^2)}}
  {1-\sigma(a)\sigma(b)},
  \label{eq:full-kernel}
\end{equation}
where $w_B(a)=\sqrt{k}\sn(2Ka/\pi)$ and $w_C(a)=\sqrt{k}\sn(4Ka/\pi)$.

The spectral properties of the kernels determine the integrals. At this stage,
however, these kernels do not operate on the same sectors. Under
$a\mapsto\pi-a$, $\sigma(a)$ and $w_B(a)$ are even, whereas $w_C(a)$ is odd.
Thus the kernels have different parity: $\widehat\rho_B$ acts on the even
sector and $\widehat\rho_C$ on the odd sector, each annihilating the opposite
sector. This is not necessarily a problem, but it is more convenient to fold
these sectors onto a common half-interval.

Put $s=2Ka/\pi$, $t=2Kb/\pi$, and $\Delta(s)=1-k^2\sn^4s$. On $0<s,t<K$, define
\begin{equation}
  \begin{aligned}
    \rho_B(s,t)&=\frac{k\sn s\sn t\sqrt{\Delta(s)\Delta(t)}}
    {1-k^2\sn^2s\sn^2t},\\
    \rho_C(s,t)&=\frac{k\sn(2s)\sn(2t)\sqrt{\Delta(s)\Delta(t)}}
    {1-k^2\sn^2s\sn^2t}.
  \end{aligned}
  \label{eq:half-kernel}
\end{equation}
Reflection of one integration variable changes a row and the corresponding
column by the same sign, leaving the determinant unchanged. Consequently,
we can compress the integrals as follows
\begin{equation}
  \begin{aligned}
    \mathcal I_G(q,1;n)
    =\left(\frac{C_0(Q,q)}2\right)^n q^{-n^2/2}Z_G(n), \qquad
    Z_G(n)=\frac1{n!}\int_{(0,K)^n}
    \prod_{i=1}^n\frac{\dd s_i}{K}\,
    \det_{i,j}\!\left[\rho_G(s_i,s_j)\right].
  \end{aligned}
  \label{eq:coefficient}
\end{equation}
This is the Fermi-gas determinant
representation~\cite{Bourdier:2015wda,Du:2023pbn}, which is fully determined by
the spectral traces $\operatorname{Tr}(\rho_G^\ell)$. Both kernels act on the
common Hilbert space $\HH=L^2((0,K),\dd s/K)$, so we can focus on these
operators.

\subsection{Kernel factorisation}

Our goal is to show that the non-zero eigenvalues of the two operators $\rho_B$
and $\rho_C$ agree despite their clear differences. Distinct operators with
identical non-zero spectra are central to supersymmetric quantum
mechanics~\cite{CooperKhareSukhatme:1995}. This motivates the following
hypothetical factorisation\footnote{
  See also \cite{Hatsuda:2012hm} for an analogous factorisation relating the
  ABJM density operator to a Hankel matrix.}:
\begin{equation}
  \rho_B=\mathsf A^\dagger\mathsf A,\qquad
  \rho_C=\mathsf A\mathsf A^\dagger.
  \label{eq:decomp}
\end{equation}

These have the same non-zero eigenvalues with multiplicities. The challenge is
to construct the operator $\mathsf A$.

As a first step, the Jacobi addition formula gives a useful common presentation
of the kernels $\rho_G(s,t)=\langle s|\rho_G|t\rangle$:
\begin{equation}
  \rho_G(s,t)=m_G(s)C(s,t)m_G(t),\qquad
  C(s,t)=\cn(s-t)-\cn(s+t),
  \label{eq:common-kernel}
\end{equation}
where
\begin{equation}
  m_B(s)=\sqrt{\frac{k}{2}}\frac{\sqrt{\Delta(s)}}{\dn s},
  \qquad m_C(s)=\sqrt{2k}\frac{\cn s}{\sqrt{\Delta(s)}}.
\end{equation}
The common function has an explicit spectral decomposition derived from the
Fourier expansion of $\cn$~\cite[\S22.11]{DLMF}:
\begin{equation}
  \begin{gathered}
    C(s,t)=\sum_{r=0}^{\infty}\gamma_r\hat{e}_r(s)\hat{e}_r(t),\qquad
    \hat{e}_r(s)=\sqrt{2}\sin(\omega_rs),\qquad
    \gamma_r=\frac{\pi}{kK}\sech(2\xi_r),
  \end{gathered}
  \label{eq:convolution-series}
\end{equation}
with
\begin{equation}
  \omega_r=\frac{(2r+1)\pi}{2K},\qquad
  \xi_r=\frac{K'\omega_r}{2}.
\end{equation}
The $\hat{e}_r$ form an orthonormal basis of $\HH$.

How do we connect this to~\eqref{eq:decomp}? First focus on the $B_n$ case. We
can absorb the one-body factors into a new set of ``basis states'' $\lvert
b_r\rangle$, where $\langle s|b_r\rangle=m_B(s)\hat{e}_r(s)$.
For the $C_n$ case, similarly introduce $\lvert c_r\rangle$ satisfying
$\langle s|c_r\rangle=m_C(s)\hat{e}_r(s)$. Loosely speaking, it would be
desirable to work in a basis that roughly diagonalises $\mathsf A$. We therefore posit
that there exists auxiliary orthonormal bases $\lvert\hat{\Phi}_r\rangle$ and
$\lvert\hat{\Psi}_r\rangle$
\footnote{One might suggest using $\lvert \hat{e}_r\rangle$ instead, which
  would satisfy $\rho_B=\mathsf A_B^\dagger \mathsf A_B$. This would not, however, satisfy
  $\rho_C=\mathsf A_B \mathsf A_B^\dagger$.}
that satisfy
\begin{equation}
  \mathsf A_B=\sum_{r\geq0}\sqrt{\gamma_r}\,
  \lvert\hat{\Phi}_r\rangle\langle b_r\rvert, \quad
  \mathsf A_C=\sum_{r\geq0}\sqrt{\gamma_r}\,
  \lvert c_r\rangle\langle\hat{\Psi}_r\rvert.
\end{equation}
These ans\"atze give $\mathsf A_B^\dagger\mathsf A_B=\rho_B$ and $\mathsf
  A_C\mathsf A_C^\dagger=\rho_C$. To obtain~\eqref{eq:decomp} we require
$\mathsf A_B=\mathsf A_C$. Assuming these bases are complete, it is enough to
match all matrix elements:
\begin{equation}
  \langle\hat{\Phi}_r|\mathsf A_B|\hat{\Psi}_\ell\rangle
  =\langle\hat{\Phi}_r|\mathsf A_C|\hat{\Psi}_\ell\rangle.
  \label{eq:matrix-elements}
\end{equation}
To make progress we return to the prior position-basis with
$\hat{\Phi}_r(s)=\langle s|\hat{\Phi}_r\rangle$,
$\hat{\Psi}_r(s)=\langle s|\hat{\Psi}_r\rangle$, and
$A_G(s,t)=\langle s|\mathsf A_G|t\rangle$. The two expansions then read
\begin{equation}
  \begin{aligned}
    A_B(s,t)&=m_B(t)\sum_{r\geq0}\sqrt{\gamma_r}\,
    \hat{\Phi}_r(s)\hat{e}_r(t),\\
    A_C(s,t)&=m_C(s)\sum_{r\geq0}\sqrt{\gamma_r}\,
    \hat{e}_r(s)\hat{\Psi}_r(t).
  \end{aligned}
  \label{eq:darboux-series}
\end{equation}
Establishing appropriate functional forms for $\hat{\Phi}_r$ and
$\hat{\Psi}_r$ is a nontrivial exercise.
In appendix~\ref{app:mode-motivation} we motivate the following construction.
Here we simply state the expressions:
\begin{equation}
  \begin{aligned}
    \hat{\Phi}_r(s)&=\frac{2\bigl[\cosh\xi_r\cos\theta(s)\sin(\omega_rs)
        +\sinh\xi_r\sin\theta(s)\cos(\omega_rs)\bigr]}
    {\sqrt{\cosh(2\xi_r)}},\\
    \hat{\Psi}_r(s)&=\frac{2\bigl[\cosh\xi_r\cos\varphi(s)\sin(\omega_rs)
        -\sinh\xi_r\sin\varphi(s)\cos(\omega_rs)\bigr]}
    {\sqrt{\cosh(2\xi_r)}},
  \end{aligned}
  \label{eq:outgoing}
\end{equation}
where
\begin{equation}
  \begin{aligned}
    \cos\theta(s)&=\frac{\cn s}{\sqrt{\Delta(s)}},&
    \sin\theta(s)&=\frac{\sn s\dn s}{\sqrt{\Delta(s)}},\\
    \cos\varphi(s)&=\frac{\dn s}{\sqrt{\Delta(s)}},&
    \sin\varphi(s)&=\frac{k\sn s\cn s}{\sqrt{\Delta(s)}}.
  \end{aligned}
  \label{eq:half-angles}
\end{equation}

We then have the following.

\begin{proposition}\label{prop:modes}
  The systems $(\hat{\Phi}_r)_{r\geq0}$ and $(\hat{\Psi}_r)_{r\geq0}$ are
  complete orthonormal bases of $\HH$. The series~\eqref{eq:darboux-series}
  converge uniformly on $[0,K]^2$ and satisfy, for every $r,\ell$,
  \begin{equation}
    \begin{aligned}
      &\int_0^K\!\int_0^K
      \hat{\Phi}^*_r(s)A_B(s,t)\hat{\Psi}_\ell(t)
      \frac{\dd s\,\dd t}{K^2}=\int_0^K\!\int_0^K
      \hat{\Phi}^*_r(s)A_C(s,t)\hat{\Psi}_\ell(t)
      \frac{\dd s\,\dd t}{K^2}.
    \end{aligned}
    \label{eq:bounded-coefficients}
  \end{equation}
\end{proposition}
\begin{proof}
  See appendix~\ref{app:darboux}.
\end{proof}

Equation~\eqref{eq:bounded-coefficients} is equivalent to
\eqref{eq:matrix-elements}. Thus we have found a sufficient set of basis
functions and consequently constructed the desired operator $\mathsf A:=\mathsf
A_B=\mathsf A_C$. Equation~\eqref{eq:decomp} holds and thus the two indices are
identical.
